\section*{第 \ref{ch:g10:numbers} 章——数与数集}

\begin{solution}{exo:g10:numbers:1}
$\frac{15}{3} = 5 \in \N$. \quad $-7 \in \Z$. \quad
$\frac{22}{7} \in \Q$（它不是\omterm{def:g10:numbers:sets}{整数}：$22 = 7 \times 3 + 1$）。
\quad $\sqrt 9 = 3 \in \N$。\quad $\sqrt{10} \in \R$（\omterm{ex:g10:numbers:classify}{无理数}，因为
$10$ 不是\omterm{def:g10:numbers:sets}{有理数}的平方——在本阶段不加证明地承认，思路同
\cref{thm:g10:numbers:sqrt2}）。\quad $-2.4 = -\frac{24}{10} \in \Q$。
\quad $0 \in \N$.
\end{solution}

\begin{solution}{exo:g10:numbers:2}
(a) $\intco{-2}{5}$：在 $-2$ 处画实心点，在 $5$ 处画空心点。
(b) $\intoo{3}{+\infty}$：在 $3$ 处画空心点，向右涂阴影。
(c) $\intoc{-\infty}{-1}$：从左侧涂阴影至
$-1$ 处的实心点。
(d) “$x$ 到 $2$ 的距离至多为 $3$” 即 $\abs{x - 2} \leq 3$，
即 $x \in \intcc{-1}{5}$：在 $-1$ 与 $5$ 处画实心点。
\end{solution}

\begin{solution}{exo:g10:numbers:3}
(a) 两个\omterm{def:g10:numbers:interval}{区间}在 $0$ 与 $2$ 之间有公共部分：
$I \cap J = \intcc{0}{2}$（$0 \in J$ 且 $2 \in I$，两数也都属于
另一集合），且 $I \cup J = \intco{-3}{4}$。

(b) $I \cap J$ 是由 $x$ 组成的集合，满足 $-2 \leq x$ 且 $x < 1$：
$\intco{-2}{1}$。\omterm{def:g10:numbers:interunion}{并集}为全体\omterm{def:g10:numbers:sets}{实数}：$I \cup J = \R$。
\end{solution}

\begin{solution}{exo:g10:numbers:4}
$\abs{-6} = 6$; \quad $\abs{4 - 9} = \abs{-5} = 5$; \quad
$\abs{-3 - 5} = \abs{-8} = 8$; \quad
$\sqrt 2 > 1$，故 $\abs{\sqrt2 - 1} = \sqrt2 - 1$；\quad
$1 - \sqrt2 < 0$，故 $\abs{1 - \sqrt2} = \sqrt2 - 1$ 亦成立（一个数与其相反数
的\omterm{def:g10:numbers:abs}{绝对值}相等）。
\end{solution}

\begin{solution}{exo:g10:numbers:5}
$\abs{x} = 5$：$5$ 是到 $0$ 的距离，故 $x = 5$ 或 $x = -5$。

$\abs{x - 1} = 3$：$3$ 是到 $1$ 的距离，故 $x = 4$ 或 $x = -2$。

$\abs{x - 4} \leq 1$：$1$ 为到 $4$ 的距离的上界，故
$x \in \intcc{3}{5}$。

$\abs{x + 2} < 3$：$3$ 为到 $-2$ 的距离的严格上界，故
$x \in \intoo{-5}{1}$。
\end{solution}

\begin{solution}{exo:g10:numbers:6}
每个\omterm{def:g10:numbers:interval}{区间}都由其中心 $a$（两端点的中点）和
半径 $r$（长度的一半）来描述。

$\intcc{2}{8}$：$a = 5$，$r = 3$，故 $\abs{x - 5} \leq 3$。

$\intoo{-5}{1}$：$a = -2$，$r = 3$，端点为开的，故 $\abs{x + 2} < 3$。

$\intcc{-7}{-3}$：$a = -5$，$r = 2$，故 $\abs{x + 5} \leq 2$。
\end{solution}

\begin{solution}{exo:g10:numbers:7}
设 $x = 0.272727\dots$，则 $100x = 27.2727\dots$，两式相减：
\[
100x - x = 27.2727\dots - 0.2727\dots = 27,
\]
故 $99x = 27$，$x = \frac{27}{99} = \frac{3}{11}$，是两\omterm{def:g10:numbers:sets}{整数}之商，
因而 $x$ 是\omterm{def:g10:numbers:sets}{有理数}。
\end{solution}

\begin{solution}{exo:g10:numbers:8}
\emph{1. 正确：}\omterm{def:g10:numbers:sets}{整数}（包括正\omterm{def:g10:numbers:sets}{整数}与负\omterm{def:g10:numbers:sets}{整数}）相加
总得到一个\omterm{def:g10:numbers:sets}{整数}。

\emph{2. 错误：}$\frac{1}{2}$ 是\omterm{def:g10:numbers:sets}{整数} $1$ 与 $2$ 的商，
但它不是\omterm{def:g10:numbers:sets}{整数}。

\emph{3. 正确：}$\frac pq + \frac{p'}{q'} = \frac{pq' + p'q}{qq'}$ 仍然是
两\omterm{def:g10:numbers:sets}{整数}之商（分母不为零）。

\emph{4. 正确：}假设 $r$ 是\omterm{def:g10:numbers:sets}{有理数}，$t$ 是\omterm{ex:g10:numbers:classify}{无理数}，而 $r + t = s$
是\omterm{def:g10:numbers:sets}{有理数}。于是 $t = s - r$ 为两个\omterm{def:g10:numbers:sets}{有理数}之差，
因而为\omterm{def:g10:numbers:sets}{有理数}（由第 3 小题，取 $-r$ 应用）——矛盾。故 $r + t$ 是
\omterm{ex:g10:numbers:classify}{无理数}。
\end{solution}

\begin{solution}{exo:g10:numbers:9}
将 $1.414 \leq \sqrt2 \leq 1.415$ 乘以 $2 > 0$：
$2.828 \leq 2\sqrt2 \leq 2.830$.

加上 $3$：$4.414 \leq \sqrt2 + 3 \leq 4.415$。

乘以 $-1 < 0$ 会使不等号方向改变：
$-1.415 \leq -\sqrt2 \leq -1.414$.
\end{solution}

\begin{solution}{exo:g10:numbers:10}
先证明一个辅助事实。对 $p$ 作除以 $3$ 的带余除法，余数为 $0$、$1$ 或
$2$，即 $p = 3k$、$p = 3k+1$ 或 $p = 3k+2$。两边平方：
\[
(3k)^2 = 3(3k^2), \quad
(3k+1)^2 = 3(3k^2 + 2k) + 1, \quad
(3k+2)^2 = 3(3k^2 + 4k + 1) + 1 .
\]
只有第一式是 $3$ 的倍数，因此若 $p^2$ 是 $3$ 的倍数，
则 $p$ 也是。

现设 $\sqrt3 = \frac pq$ 为最简分数。两边平方得
$p^2 = 3q^2$，故 $p^2$ 是 $3$ 的倍数，从而 $p = 3k$。于是
$9k^2 = 3q^2$，$q^2 = 3k^2$，因此 $q$ 也是 $3$ 的倍数——但
则分数 $\frac pq$ 就不是最简分数，与假设矛盾。故
$\sqrt3$ 是\omterm{ex:g10:numbers:classify}{无理数}。
\end{solution}

\begin{solution}{pb:g10:numbers:1}
\textbf{1.} $-7 \in \Z$; \ $\frac{13}{4} \in \Q$; \
$\sqrt{16} = 4 \in \N$; \ $0.121212\ldots = \frac{12}{99} =
\frac{4}{33} \in \Q$（初中分册中的循环小数周末问题）；\
$\sqrt8 = 2\sqrt2$ 是\omterm{ex:g10:numbers:classify}{无理数}：最小的数集为 $\R$；\ $\pi$：
$\R$.

\textbf{2.} $\frac pq + \frac rs = \frac{ps + qr}{qs}$ 及
$\frac pq \times \frac rs = \frac{pr}{qs}$：均为两\omterm{def:g10:numbers:sets}{整数}之
商（分母非零）——\omterm{def:g10:numbers:sets}{有理数}。

\textbf{3.} (a) 若 $r$ 为\omterm{def:g10:numbers:sets}{有理数}，$x$ 为\omterm{ex:g10:numbers:classify}{无理数}，且
$r + x = q$ 为\omterm{def:g10:numbers:sets}{有理数}，则 $x = q - r$ 为两个\omterm{def:g10:numbers:sets}{有理数}之差，
故为\omterm{def:g10:numbers:sets}{有理数}（第 2 题）：
矛盾。(b) 若 $r \neq 0$ 且 $r x = q$ 为\omterm{def:g10:numbers:sets}{有理数}，则
$x = \frac qr$ 为\omterm{def:g10:numbers:sets}{有理数}：矛盾。

\textbf{4.} $\sqrt2$ 与 $-\sqrt2$ 均为\omterm{ex:g10:numbers:classify}{无理数}，其和为 $0$；
$\sqrt2$ 与 $\sqrt2$ 之积为 $2$。有理的结果来自无理的
成分：\omterm{ex:g10:numbers:classify}{无理数}并不构成稳定的“王国”。

\textbf{5.} $\sqrt8 = 2\sqrt2$，故
$\sqrt2 + \sqrt8 = 3\sqrt2$：为\omterm{ex:g10:numbers:classify}{无理数}（第 3 题 (b)，乘数
$3$）。又 $\sqrt2 \times \sqrt8 = \sqrt{16} = 4 \in \N$：
两个\omterm{ex:g10:numbers:classify}{无理数}之积竟落入
所有数中最小的“王国”。

\textbf{6.} $3.475$，再取 $3.471$（或 $3.4701$，$3.47001$，
\dots）：每添一位小数，便产生新的\omterm{def:g10:numbers:sets}{有理数}，
永远介于两者之间。

\textbf{7.} $10^{-n}$ 的各倍数沿数轴排列，步长
为 $10^{-n} < g$。第一个严格大于
$a$ 的倍数——它必然存在，因为倍数终将超过
$a$——至多超出 $a$ 一步，故位于
$a + g = b$ 之前，严格介于 $a$ 与 $b$ 之间。格点即为
小数，从而为\omterm{def:g10:numbers:sets}{有理数}：长度为正的任一\omterm{def:g10:numbers:interval}{区间}
都含有一个这样的点。

\textbf{8.} 在 $a$ 与第 7 题中得到的\omterm{def:g10:numbers:sets}{有理数} $r_1$ 之间，
仍由第 7 题知存在\omterm{def:g10:numbers:sets}{有理数} $r_2$；在
$a$ 与 $r_2$ 之间再取\omterm{def:g10:numbers:sets}{有理数} $r_3$；依此类推，得无穷多个，
两两互异，且都在原\omterm{def:g10:numbers:interval}{区间}之内。

\textbf{9.} $r + \frac{\sqrt2}{10^k}$ 是\omterm{def:g10:numbers:sets}{有理数}与
\omterm{ex:g10:numbers:classify}{无理数}之和（$\frac{\sqrt2}{10^k}$ 由第 3 题 (b) 知为\omterm{ex:g10:numbers:classify}{无理数}），
故此和为\omterm{ex:g10:numbers:classify}{无理数}。由于
$\frac{\sqrt2}{10^k} < \frac{2}{10^k}$ 可任意小，
故当 $k$ 充分大时，$r + \frac{\sqrt2}{10^k}$ 仍位于
$b$ 之前：\omterm{def:g10:numbers:interval}{区间}内的一个\omterm{ex:g10:numbers:classify}{无理数}——改变
$k$ 即得无穷多个。

\textbf{10.} 正\omterm{def:g10:numbers:sets}{实数}中无最小者：若 $s > 0$ 为最小者，则
\omterm{def:g10:numbers:interval}{区间} $(0, s)$ 仍含一个\omterm{def:g10:numbers:sets}{有理数}（第 7 题），
该\omterm{def:g10:numbers:sets}{有理数}为正且小于 $s$。$3$ 之后不存在紧邻的数：任一
候选数 $c > 3$ 都使\omterm{def:g10:numbers:interval}{区间} $(3, c)$ 不
为空——即初中分册中“无下一个数”周末问题的游戏，如今成为关于
$\R$.

\textbf{11.} $\abs{x - 5} = 2$：$x = 3$ 或 $x = 7$。
$\abs{x - 5} < 2$: $x \in \intoo{3}{7}$.
$\abs{x + 1} \geq 3$：到 $-1$ 的距离至少为 $3$：
$x \in \intoc{-\infty}{-4} \cup \intco{2}{+\infty}$.

\textbf{12.} $\abs{d - 80} \leq 0.05$，即\
$d \in \intcc{79.95}{80.05}$。活塞在 $79.97$ 处合格；
在 $80.06$ 处的活塞不合格，超差百分之一毫米。

\textbf{13.} $(3, -5)$：$\abs{-2} = 2 \leq 8$。$(-2, -7)$：
$\abs{-9} = 9 = 2 + 7$：取等号。同号时，$\abs{a + b}$ 为
两距离之和，等于 $\abs a + \abs b$。异号时，
折返而行，$\abs{a + b}$ 为
两距离之\emph{差}，严格小于其
和（除非其中之一为 $0$）。取等号当且仅当 $a$ 与
$b$ 同号或其中之一为零。

\textbf{14.} 解为到 $2$
与 $-4$ 等距离的点：中点 $x = -1$。它是
初中分册中“两面镜子”周末问题的垂直平分线，压缩到一维后即为
一个点。

\textbf{15.} $\sqrt{x^2} = \abs{x}$，而非 $x$：当 $x = -3$ 时，
$\sqrt{9} = 3 = \abs{-3} \neq -3$。于是 $x^2 < 9$ 化为
$\sqrt{x^2} < 3$，即\ $\abs x < 3$：
$x \in \intoo{-3}{3}$.

\textbf{16.} 测量只能断定长度位于
\omterm{def:g10:numbers:interval}{区间} $\intcc{1.41420}{1.41422}$ 内——而由
第 7 题与第 9 题可知，\emph{该\omterm{def:g10:numbers:interval}{区间}中含有无穷多个\omterm{def:g10:numbers:sets}{有理数}，
也含有无穷多个\omterm{ex:g10:numbers:classify}{无理数}}。对于任何将来的
容差，无论多么小，同样成立：测量永远无法
区分这两类数；只有关于精确长度的证明
——如单位正方形对角线那样的——
（初中分册中的无理性周末问题）才能做到。

\textbf{17.} 它们逼近 $\sqrt2$，误差分别小于
$10^{-1}, 10^{-2}, 10^{-3}, \dots$：\omterm{def:g10:numbers:sets}{有理数}在每一尺度上都紧贴
$\sqrt2$。一般地，每个\omterm{def:g10:numbers:sets}{实数}都可
被\omterm{def:g10:numbers:sets}{有理数}任意逼近（即其十进小数
截断）——这正是从目标一侧看到的稠密性，
再一次体现。

\textbf{18.} $a + b \in \intcc{2.4 + 1.1}{2.6 + 1.3} =
\intcc{3.5}{3.9}$：不确定度为 $\pm 0.2$（误差相加）。
$a \times b \in \intcc{2.4 \times 1.1}{2.6 \times 1.3} =
\intcc{2.64}{3.38}$：不确定度约为 $\pm 0.37$，位于
$3$ 附近——乘积的相对误差相加，因此精度
下降得更快。

\textbf{19.} 写出真值 $a = 2.5 + e_1$，
$b = 1.2 + e_2$，其中 $\abs{e_1}, \abs{e_2} \leq 0.1$。于是
和的误差为 $\abs{e_1 + e_2} \leq \abs{e_1} +
\abs{e_2} \leq 0.2$：工程师的这条法则\emph{正是}第 13 题的
三角不等式。

\textbf{20.} \omterm{def:g10:numbers:sets}{实数}轴既无空隙，也无相邻之点：任一数之后
都没有“下一个”数（第 10 题）。其点集分为
\omterm{def:g10:numbers:sets}{有理数}——恰为有限小数或循环小数
（初中分册的循环小数周末问题）——与\omterm{ex:g10:numbers:classify}{无理数}两类，
两类数交织得如此紧密，以至于每一
\omterm{def:g10:numbers:interval}{区间}内都各含无穷多个（第 8 题与第 9 题），
这正是为何只有证明，而非测量，才能将二者区分
（第 16 题）。最后的惊奇，作为承诺留待下文：
\omterm{ex:g10:numbers:classify}{无理数}的个数多于\omterm{def:g10:numbers:sets}{有理数}——并非一个一个地数，
而是在比较无穷集合大小的精确意义上，
这需要一套在本科分册中建立的理论。
\end{solution}
